\documentclass[11pt]{article}
\usepackage[margin=1in]{geometry}
\usepackage{amsmath,amssymb,hyperref}
\newcommand{\Fp}{\mathbb F_p}
\title{Joint recovery of the exceptional multiplicative contribution}
\author{Project research note}
\date{18 November 2026}
\begin{document}
\maketitle

In the paired-slot model the exceptional product can be computed exactly
from the seeds in $1+\pi(\lfloor\sqrt p\rfloor)$ ordinary count queries per
component, independent of the number of distinct exceptional weights.
Each query has an arithmetic graph with $O(mb+b+L)$ Boolean edges,
where $b=\lceil\log_2p\rceil$ and $L$ is the bit length of the public
integer cover. For $L=O(b)$ it meets the project's square-root degree
and polylogarithmic parallel-edge bounds at tuned arity. There is also an adaptive alternative that
can certify simultaneous cancellation. Neither gives a better general
end-to-end exponent.

\section{Which product is recovered}
For a nonzero local field weight $w$ write $h=\min(w,p-w)$ using its
canonical representative. For a public factored integer $F>0$ set
$r=h/\gcd(h,F)$. The desired component correction is
\[
                         E_F=\prod_{\text{oriented occurrences}} r \pmod p .
\]
Covered weights have $r=1$. The sign of $w$ and the prime powers in
$\gcd(h,F)$ belong to the previously computed covered factor. Use
the terminal weights and prefix chord-scale weights of
\texttt{target.tex}, including their identity, tangent and target
patches. The final correction multiplies the terminal result by the
prefix result for $j$ raised to $\prod_{i>j}2n_i$. An empty slot gives
correction one. Inverse dyadic windows give the slot sizes $n_i$ in
$O(\sum_i\tau_i\,\mathrm{poly}(b))$ time. No half-list is needed.

\section{Two short integers for every cofactor}
Set $B=\lfloor\sqrt p\rfloor$, $x=\min(r,r^{-1}\bmod p)$.
Let $u/v<x/p<u'/v'$ be the consecutive fractions surrounding $x/p$
in the Farey sequence of order $B$. They exist since $x/p$ is reduced
with denominator $p>B$. Their certificate consists of the two strict
inequalities and
\[
 0\le u\le v,\quad 0\le u'\le v',\quad 1\le v,v'\le B,\quad
 u'v-uv'=1,\quad v+v'>B .
\]
The determinant forces reduced fractions, and a fraction between
them must have denominator at least $v+v'$. Indeed, write its
numerator--denominator vector in the unimodular basis of the two
endpoints. Both coefficients are positive integers. The standard
Farey neighbour property gives existence, and disjoint interiors give
uniqueness.\footnote{See the primary paper
\href{https://math.colgate.edu/~integers/d7/d7.pdf}{The relative size of consecutive odd denominators in Farey series}.}
The scalar verifier finds the endpoints by batched mediant steps. A
count checks the certificate directly and has no mediant loop.

Write $a=xv-pu>0$ and $a'=pu'-xv'>0$. Then
$av'+a'v=p$. Consequently
\[
    \min(a,a')\le\left\lfloor p/(B+1)\right\rfloor\le B .
\]
On $a\le a'$ choose the positive ratio $a/v$. Otherwise choose the
negative ratio $-a'/v'$. It equals $x$ modulo $p$. Swap numerator and
denominator if $r\ne x$. This gives a canonical triple
\[
        r=(-1)^\epsilon A/D\pmod p,\qquad 1\le A,D\le B .
\]
Inverting a residue that is not self-inverse swaps its two labels and
preserves the sign. The convention $r=1$ gives $+1/1$.

\paragraph{Unique arithmetic witnesses.}
First obtain $t=F\bmod h$ by the public-limb recurrence in
\texttt{unit\_recovery.tex}. Check $h=gr,\ t=gu$ with $g,r>0$ and
$0\le u<r$. For $r>1$, require $uz=1+kr$ with $0<z<r$ and
$0\le k<r$. This certifies $\gcd(u,r)=1$ and fixes $z,k$.
For $r=1$ force $u=z=k=0$. These conditions force $g=\gcd(h,F)$.
The canonical field inverse, Farey inequalities, determinant,
errors, sign and swap each use a constant number of integer or
modular products and comparisons on $O(b)$-bit words. Together with
the remainder computation they add $O(b+L)$ edges. Each product uses
the existing row-distributed multiplication check. Complete words
can use the original owners and constant-length equality relays.
The fixed number of copies per stage preserves the required degree
and parallel multiplicity bounds.

\section{Count each prime's contribution across all exceptional weights}
For every prime $q\le B$, let $J_q=\lfloor\log_q B\rfloor$. For one
component define
\[
  e_q=\sum_{j=1}^{J_q}\bigl(\#\{q^j\mid A\}-\#\{q^j\mid D\}\bigr).
\]
The answer is $E_F=(-1)^{C_-}\prod_{q\le B}q^{e_q}\bmod p$, where
$C_-=\sum\epsilon$. Negative exponents use modular inverses.
Every query depends on the fixed prime or prime power and the seeds.
There is no search for individual exceptional magnitudes.

The version using the original guarded path requests one nonnegative
pin per divisibility test and one for the sign. Each has unique
auxiliaries, so the earlier box rounding and sign-symmetry division
apply verbatim. Its number of predicate types per component is
$1+2\sum_{q\le B}J_q$, before dyadic box splitting.

For the version using one query per prime, let $M$ be the exact
oriented population. With $j$ slots, $M=\prod_{i\le j}(2n_i)/2$ in
the terminal component and $\prod_{i\le j}(2n_i)/4$ in a prefix
component, where $j\ge2$. A witness chooses a public-power index
$s\in[1,J_q]$ and one of two branches. Accept $q^s\mid A$ on the
first branch and $q^s\nmid D$ on the second. The nonnegative count
is exactly $J_q M+e_q$. Subtract the known offset in ordinary
integer arithmetic. The $O(\log b)$ index bits select among at most
$b$ public powers of $q$ computed in $\mathrm{poly}(b)$ time per
prime. The existing row owners check their own bits of the selected
word locally from the copied index. This costs
$O(b+\sqrt b\,\mathrm{polylog}(b))$ edges, and canonical division
supplies the quotient and remainder.

Use ordinary oriented counts for these deliberate witness
multiplicities. In particular one must not send their weighted raw
total through the old single-witness rounding formula on the old
boxes. An explicit ordinary graph uses canonical affine coordinates
plus an infinity flag with fixed dummy coordinates. For each
addition, disjoint branches handle infinity, inverse pairs, doubling
and distinct-abscissa addition. Active slopes and modular quotients
are unique. Set unused auxiliary ports to zero on every branch.
Constrain every input to the curve and its hash window. Select the
canonical lower ordinate in the first slot and additionally in the
last slot of a prefix component. This realizes the oriented count
without an offset. Each addition has a fixed number of field
equations, so compiler copies cost $O(mb)$ edges. This argument
does not transfer the optimized coefficient 36 to this graph.

Primes can be streamed by scanning to $B$ with deterministic
polynomial-bit primality tests. Preprocessing costs
$O(B\,\mathrm{poly}(b))$ time and polynomial working storage.
The number of primes has $\log_2(1+\pi(B))=\tfrac12 b+O(\log b)$.
An elementary lower estimate suffices here. For $n=\lfloor B/2\rfloor$,
$\binom{2n}{n}\ge 2^{2n}/(2n+1)$ divides $\mathrm{lcm}(1,\dots,B)$
because each prime valuation of that binomial is at most
$\lfloor\log_q(2n)\rfloor$. The lcm has log size at most
$\pi(B)\log_2 B$. Thus the fixed sweep alone has a linear exponent
in $b$. For $L=O(b)$ the guarded implementation has conditional
one-unit time $2^{b/2+o(b)}$ at the tuned widths, including the
subboxes, polynomially many components and preprocessing. This is a
bound for computing the correction without half-lists.

\section{When numerator and denominator counts cancel}
For a label $d\le B$ put $c_d=N_A(d)-N_D(d)$. For an interval $I$
compute its imbalance by exact ordered-pair counts
\[
  H_I=\#AA_I-2\#AD_I+\#DD_I=\sum_{d\in I}c_d^2 .
\]
For example $AD_I$ counts independent oriented occurrences whose
first numerator equals the second denominator and lies in $I$.
Include the diagonal pairs. Each of the three counts uses two
independent copies of the ordinary graph above with a single equality
and interval predicate. Two single-copy counts give
$S_I=\sum_{d\in I}c_d$. All counts are integers. A zero $H_I$ proves
the entire interval contributes one. Bisect the other intervals.
At a singleton multiply by $d^{S_I}$. With $K$ nonzero labels this
takes $O(1+K\log B)$ interval tests, capped by $2B-1$ for the binary
tree. A single root test suffices when all labels cancel. The sign
is still counted separately. The two-copy graph has a larger
arithmetic constant than a one-copy count.

\paragraph{A seed family with many surviving labels.}
Choose a dyadic $W$ of order $B/4$ with $W\le B/2$. Partition the
hash modulus $2^b$ into intervals of width $W$. Since the group has
$\Theta(p)$ rational points in paired signs, some interval contains
$\Omega(W)$ distinct admissible abscissas. Take that interval as a
one-slot affine hash window with multiplier one and choose a target
in the same interval. For $F=1$, all unpatched magnitudes are positive
distances $h<W$, and at most two abscissas give any one distance.
For $1<h\le B/2$ the inverse representative exceeds $h$ and the
bracket for $h/p$ is $(0/1,1/B)$. Its chosen numerator is $h$.
Thus every one of $\Omega(W)$ distinct such labels is $h/1$, with
no denominator to cancel it. The imbalance search has $\Omega(B)$
occupied leaves in this allowable family. This example concerns a
one-slot seed and does not assert a lower bound for the tuned
independently sampled multi-slot distribution.

There is also a limit on uniform two-label representations themselves.
There are at most $2Y^2$ signed ratios with both labels at most $Y$.
With one full slot and $F=1$, the at-most-four-preimage argument of
\texttt{unit\_recovery.tex} realizes $\Omega(p)$ distinct rough values.
Any scheme of this two-label form covering that entire input family
requires $Y=\Omega(\sqrt p)$. This counts available representations.
It permits methods that combine a large label domain without visiting it.

\section{Finite evidence and end-to-end accounting}
Run \texttt{python -B research/single\_target/verify\_exception\_product.py}.
The verifier compares 2,864 brackets with independently sorted
fractions and checks the ratio invariant on fields through 127 bits.
Across 22 seed/cover cases, the prime-power, weighted-prime and
collision reductions recover the literal correction, including
the suffix exponents. Multiplication by the independently computed
covered factor gives the literal normalized polynomial unit.
On the earlier $\mathbb F_{211}$ fixture with $F=12$, 78 distinct
terminal exceptions give correction 100. This takes seven ordinary
weighted count types, or 19 guarded pin types before box splitting.
The correction including the prefix is 132 and the full unit is 93.
On the two equal-moment fixtures the corrections are 144 and 202,
giving the respective full units 104 and 200.

On twelve independently drawn four-slot hash fixtures over
$\mathbb F_{65521}$ (widths $8,16,16,16$, target $(0,1)$ and $F=12$)
the terminal components had 5,975 to 21,572 distinct exceptional
magnitudes. All products matched the 55 prime/sign coordinates.
Between 239 and 254 of the 255 integer labels had nonzero imbalance.
The finite probe streams signed tuples and records reference
histograms to measure these statistics. As a separate control,
56 reciprocal-paired rough values cancelled in one root test.
These are arithmetic checks. They do not time the conditional
large-instance contractor.

The fixed prime sweep already costs $2^{b/2+O(\log b)}$ count types
before applying the contractor. The imbalance variant would need
quantitative control of cancellations for the tuned seed family,
as well as a bound for the paired queries. Neither route reduces the
training term or establishes faster relation extraction. The best
project solver estimate under its finite-packet and contraction
assumptions remains
\[
 \log_2\mathbb E T_{\rm solve}\le (26.2933775156+o(1))\,
                                   b/(\log_2 b)^{1/4}.
\]
The follow-up \texttt{prime\_projection.tex} combines the prime
columns, terminal component and prefix components into one integer
count, and prices its exponent certificate against both contraction
theorems. It also tests reconstruction from small auxiliary fields.
\end{document}
